\chapter{Symmetrische Komponenten}
\label{chap:07-symmetrische-komponenten}

\section{Motivation}

Die Methode der symmetrischen Komponenten ermöglicht es, ein unsymmetrisches
Dreiphasensystem in drei symmetrische Teilsysteme zu zerlegen. Sie ist ein
zentrales Verfahren zur Analyse unsymmetrischer Kurzschlüsse und Erdschlüsse.

Die Methode ist besonders leistungsfähig, wenn die Betriebsmittel und das
Netz im betrachteten Modell linear beschrieben werden können. Bei stark
nichtlinearen Betriebsmitteln oder leistungselektronischen Einspeisern muss
geprüft werden, ob das zugrunde gelegte Modell die tatsächliche Fehlerreaktion
hinreichend beschreibt.

\section{Definition der symmetrischen Komponenten}

Es sei
\[
  a=e^{j120^\circ}
  =-\frac12+j\frac{\sqrt{3}}{2}.
\]
Dann gilt
\[
  a^2=e^{j240^\circ}
  =-\frac12-j\frac{\sqrt{3}}{2}
\]
sowie
\[
  a^3=1,\qquad 1+a+a^2=0.
\]

Die Phasenspannungen \(U_a,U_b,U_c\) werden in Null-, Mit- und
Gegensystemkomponenten zerlegt:
\[
  \begin{bmatrix}
    U_0\\
    U_1\\
    U_2
  \end{bmatrix}
  =
  \frac13
  \begin{bmatrix}
    1&1&1\\
    1&a&a^2\\
    1&a^2&a
  \end{bmatrix}
  \begin{bmatrix}
    U_a\\
    U_b\\
    U_c
  \end{bmatrix}.
\]

Die Rücktransformation lautet
\[
  \begin{bmatrix}
    U_a\\
    U_b\\
    U_c
  \end{bmatrix}
  =
  \begin{bmatrix}
    1&1&1\\
    1&a^2&a\\
    1&a&a^2
  \end{bmatrix}
  \begin{bmatrix}
    U_0\\
    U_1\\
    U_2
  \end{bmatrix}.
\]

Dieselbe Transformation kann für die Ströme verwendet werden. Die
Phasenfolge und die Vorzeichenkonvention müssen dabei durchgehend
konsistent sein.

\section{Physikalische Bedeutung}

\subsection{Mitsystem}

Das Mitsystem besitzt die normale Phasenfolge \(a\)-\(b\)-\(c\).
Bei einem symmetrischen Netzbetrieb ist es das dominierende System.

\subsection{Gegensystem}

Das Gegensystem besitzt die umgekehrte Phasenfolge. Es wird insbesondere
bei unsymmetrischen Belastungen und Fehlern angeregt.

Bei rotierenden Maschinen können Gegensystemströme zusätzliche Verluste
und thermische Beanspruchungen verursachen. Die zulässige Beanspruchung
ist deshalb betriebsmittelabhängig.

\subsection{Nullsystem}

Im Nullsystem sind alle drei Phasoren gleichphasig:
\[
  U_{a0}=U_{b0}=U_{c0}=U_0.
\]

Die Nullsystemkomponente ist eng mit dem Neutralleiter- und Erdungsrückweg
verbunden. Ob ein Nullsystemstrom fließen kann und wie groß die
Nullsystemimpedanz ist, hängt von der Netz- und Betriebsmittelstruktur ab.

\section{Nullsystem und Neutralleiterstrom}

Für die Phasenströme gilt
\[
  I_a+I_b+I_c=3I_0.
\]

Der Neutralleiterstrom hängt von der gewählten Stromrichtung ab. Wird
der Neutralleiterstrom vom Sternpunkt zum Netz als positiv definiert,
kann beispielsweise
\[
  I_N=-(I_a+I_b+I_c)=-3I_0
\]
gelten.

Die Gleichung verdeutlicht, dass sich die drei Phasenströme bei einem
symmetrischen Betrieb ohne Nullsystemkomponente vektoriell zu null addieren.
Bei einem Erdschluss oder einer unsymmetrischen Belastung kann eine
Nullsystemkomponente auftreten.

\section{Sequenzimpedanzen}

Für die drei Sequenzsysteme werden getrennte Ersatzimpedanzen angesetzt:
\[
  Z_1,\qquad Z_2,\qquad Z_0.
\]

Die Mitsystemimpedanz \(Z_1\) beschreibt die Reaktion auf Mitsystemgrößen.
Die Gegensystemimpedanz \(Z_2\) beschreibt die Reaktion auf Gegensystemgrößen.
Die Nullsystemimpedanz \(Z_0\) hängt wesentlich von den Nullsystem-Rückwegen ab.

Für viele passive Betriebsmittel sind Mit- und Gegensystemimpedanzen ähnlich.
Diese Ähnlichkeit darf jedoch nicht ohne Prüfung als exakte Gleichheit
angenommen werden. Die Nullsystemimpedanz kann deutlich anders sein.

Bei Kabeln beeinflussen insbesondere metallische Schirme, Bewehrungen,
Erdungsverbindungen und die Anordnung der Leiter das Nullsystemverhalten.
Bei Transformatoren sind Wicklungsschaltung, Sternpunktzugänglichkeit,
Erdung und magnetischer Rückweg zu berücksichtigen.

\section{Sequenznetze und Fehlerbedingungen}

Die Sequenznetze werden entsprechend den Randbedingungen des betrachteten
Fehlers miteinander verbunden. Dabei werden die physikalischen
Fehlerbedingungen zunächst in Phasengrößen formuliert und anschließend
in Sequenzgrößen übersetzt.

Ein dreipoliger symmetrischer Kurzschluss beansprucht im idealisierten
symmetrischen Modell nur das Mitsystem. Ein einpoliger Erdschluss erfordert
dagegen im Allgemeinen die Mit-, Gegen- und Nullsysteme.

\section{Beispiel: Zerlegung eines unsymmetrischen Systems}

Gegeben seien die Phasenspannungen
\[
  U_a=230\angle 0^\circ\,\mathrm V,
\]
\[
  U_b=200\angle -125^\circ\,\mathrm V,
\]
\[
  U_c=240\angle 118^\circ\,\mathrm V.
\]

Die Sequenzkomponenten ergeben sich aus
\[
  U_0=\frac13(U_a+U_b+U_c),
\]
\[
  U_1=\frac13(U_a+aU_b+a^2U_c),
\]
\[
  U_2=\frac13(U_a+a^2U_b+aU_c).
\]

Da die Größen als komplexe Phasoren gegeben sind, müssen sie vor der
Addition in kartesische Form oder mit einer konsistenten komplexen
Rechenoperation überführt werden. Die Beträge allein dürfen nicht addiert
werden.

Die Rücktransformation muss anschließend wieder die ursprünglichen
Phasenspannungen ergeben. Sie ist eine nützliche Plausibilitätsprüfung
für eine manuelle oder softwaregestützte Rechnung.

\section{Grenzen der Methode}

Die symmetrischen Komponenten sind eine mathematische Transformation
der Phasengrößen. Sie machen ein unsymmetrisches Problem nicht automatisch
linear und beseitigen auch nicht die Notwendigkeit korrekter Betriebsmittelmodelle.

Besonders zu prüfen sind:
\begin{itemize}
  \item die Gültigkeit linearer Ersatzimpedanzen;
  \item die korrekte Nullsystemmodellierung;
  \item die Phasenfolge und die Vorzeichenkonvention;
  \item die Auswirkungen von Transformator-Schaltgruppen;
  \item die Fehlerimpedanz und der konkrete Erdungszustand;
  \item die Eigenschaften von Umrichtern und anderen geregelten Quellen.
\end{itemize}

\section{Übungsaufgaben}

\begin{enumerate}
  \item Leiten Sie die Beziehung \(I_a+I_b+I_c=3I_0\) aus der
        Rücktransformation her.
  \item Beschreiben Sie die physikalische Bedeutung von Mit-, Gegen-
        und Nullsystem.
  \item Warum ist die Nullsystemimpedanz eines Kabelsystems nicht
        ohne Weiteres gleich der Mitsystemimpedanz?
  \item Erläutern Sie, warum ein symmetrischer dreipoliger Kurzschluss
        im idealisierten Modell nur das Mitsystem benötigt.
\end{enumerate}